\section{Introduction}
\label{sec:introduction}

A minimum-variance portfolio of five US exchange-traded funds that
conditioned its covariance forecast on volatility states inferred in
real time earned an annualized Sharpe ratio \primarySharpeDiff{} higher
than the same portfolio without regime conditioning over 2010--2026, net
of 10 basis points in transaction costs. The 95\% confidence interval on
that difference runs from \primaryCiLower{} to \primaryCiUpper{}. The
strategy did what the theory predicts, realizing \volDifferenceMagnitude{}
lower annualized volatility, but it traded \turnoverRatio{} times as much
as its comparator, and the estimated advantage falls monotonically as costs
rise. On this evidence, regime conditioning delivers a measurable risk
reduction too small relative to sampling variation to establish a
risk-adjusted improvement once trading costs are paid.

The design was fixed before the data were touched. I published the
hypothesis, the comparator, the sample split, the cost convention and
all thirteen robustness specifications, and version-tagged them in a
public repository before estimating a single model. Four later
amendments each carry a date and a record of what was known when I made
them; three were adopted before any portfolio return existed. The
comparison reported here is the one committed to in advance, not the one
that survived a search. Preregistration, frozen data, causal tests,
automated validation and dated amendments substantially constrained
ex-post researcher discretion, though they did not eliminate it.

The choice of comparator does most of the analytical work. Regime-aware
minimum variance beats equal weight and 60/40 comfortably on realized
Sharpe, and a study could stop there and report an improvement. It
should not. Those benchmarks differ from the regime strategy in
optimization, covariance estimation and regime conditioning at the same
time, so a gap between them attributes nothing to any single ingredient.
I therefore compare against rolling Ledoit--Wolf minimum variance, which
shares every ingredient except the regime signal. The \primarySharpeDiff{}
difference measures what regime conditioning adds on top of dynamic
covariance estimation, and it is not distinguishable from zero.

Three findings temper the point estimate further. The sign reverses
under a three-state specification and when realized volatility is
dropped from the feature set. It reverses again between the pre- and
post-2020 halves of the sample. And after Holm adjustment across the
robustness family, no specification retains significance. Set against
these, a positive point estimate with an interval \ciWidthMultiple{} times
its width is best read as inconclusive about small effects rather than as evidence
of either a benefit or its absence.

This paper contributes in three ways. It tests a regime-conditioning
strategy under a preregistered design with a matched comparator, which
is uncommon in the backtesting literature and directly addresses the
multiple-testing critique. It reports the full ablation ladder, showing
where risk-adjusted gains actually arise: mostly from optimization
itself rather than from dynamic covariance or regime conditioning. And
it documents an interaction the design did not anticipate, in which
portfolio constraints moved the estimated difference more than any
modelling choice examined.

\section{Related literature}
\label{sec:literature}

Three strands bear directly on this study.

\citet{hamilton1989} introduced the regime-switching framework this
paper applies, and with it the discipline that states are latent and
inferred rather than observed. That distinction governs the empirical
design here: the trading signal is the filtered probability at the
decision date, never a smoothed path that conditions on data the
investor did not have. The volatility clustering the framework describes
was documented by \citet{engle1982} and generalized by
\citet{bollerslev1986}; this paper does not revisit it, asking instead
whether acting on it survives out-of-sample testing net of costs.

Applying regime switching to allocation has a substantial literature.
\citet{ang2002} document regime shifts in international equity
correlations and their allocation consequences, \citet{guidolin2007}
solve multivariate regime-switching allocation problems, and
\citet{ang2012} survey the field. This paper's contribution is narrower
and more sceptical than a survey of that work would suggest: it holds
three things fixed at once --- states inferred strictly in real time, a
comparator differing by the regime signal alone, and trading costs
charged on realized turnover --- and preregisters the comparison before
estimating anything. The question is not whether regimes exist or
whether conditioning on them helps in principle, but what the
conditioning is worth after those three constraints are imposed
simultaneously.

\citet{demiguel2009} provide the closest methodological precedent. They
evaluate fourteen portfolio models across seven datasets and find that
none consistently outperforms naive 1/N allocation on Sharpe ratio,
certainty-equivalent return or turnover, concluding that estimation
error more than offsets the theoretical gain from optimization. Their
framing is the model for this study: a negative out-of-sample result
becomes informative when it identifies the mechanism by which a
theoretically superior method fails. The question here is one rung
further along --- optimization is granted, and the test is whether
regime conditioning adds anything on top of it.

\citet{harvey2015} and \citet{harvey2016} supply the inferential
standard. Documenting that hundreds of factors have been tested in the
cross-sectional literature, they argue that conventional significance
thresholds no longer establish anything and propose haircuts for
reported Sharpe ratios that account for the number of trials.
Preregistering a single primary comparison and applying Holm adjustment
across an explicitly bounded robustness family is a direct response to
that critique: the number of trials is fixed and published in advance
rather than reconstructed afterward.

The portfolio problem is \citet{markowitz1952} with the mean removed.
\citet{michaud1989} sets out why sample mean-variance optimization
performs poorly out of sample, which is the reason no strategy here uses
expected-return forecasts, and \citet{clarke2011} characterize what
minimum-variance portfolios actually hold. \citet{jagannathan2003} show
that imposing constraints which are wrong as beliefs can nonetheless
reduce estimation error and improve realized risk --- a result that
bears directly on the constraint interaction reported in
Section~\ref{sec:robustness}.

On methods, the covariance estimator follows \citet{ledoit2004}, whose
shrinkage approach is designed for exactly the near-collinear setting
observed here, where the equity block carries pairwise correlations
between \equityCorrMin{} and \equityCorrMax{}. The hidden Markov model is estimated by the EM
algorithm of \citet{dempster1977}. Volatility forecast evaluation uses
QLIKE following \citet{patton2011}, who shows that rankings can be
distorted when noisy volatility proxies are evaluated under
inappropriate loss functions, with pairwise comparisons following
\citet{diebold1995}. Inference uses the stationary bootstrap of
\citet{politis1994}, chosen because portfolio returns are serially
dependent and independent resampling would understate uncertainty;
\citet{ledoit2008} develop the robust Sharpe-ratio test that motivates
bootstrapping the difference rather than comparing ratios directly.
Standard errors on mean differences use \citet{newey1987}, and
multiplicity is handled by \citet{holm1979}. \citet{bailey2014} give
the sharpest statement of why backtested performance requires this kind
of discipline.
